\section{Event products and elementary density constructions}
\label{sec:annex-p15-density-and-product-exercises}

\begin{annexitem}{annex:p15.event-product-criterion}{thm:p15-event-product-criterion}{Both Borel--Cantelli directions}
\begin{proof}
On an independent uniform sequence space let \(A_m=\{U_m\leq p_m\}\).
Then \(P(A_m)=p_m\), the events are independent, and continuity
of probability on decreasing events gives
\[
             P\!\left(\bigcap_{m\geq1}A_m^c\right)
                 =\lim_n\prod_{m=1}^n(1-p_m).
\]
If \(\sum_mp_m=\infty\), the second Borel--Cantelli lemma says
infinitely many \(A_m\) occur almost surely. The displayed
no-occurrence event therefore has probability zero.

If the sum is finite, the first Borel--Cantelli lemma says only
finitely many \(A_m\) occur almost surely. Write
\(T_N=\bigcap_{m>N}A_m^c\).
Then \(\bigcup_NT_N\) has probability one, so some \(T_N\)
has positive probability. Independence of the finite initial
coordinate block from the remaining coordinate sigma-field gives
\[
 P\!\left(\bigcap_{m\geq1}A_m^c\right)
           =\left(\prod_{m=1}^N(1-p_m)\right)P(T_N)>0.
\]
Each initial factor is positive by \(p_m<1\). This proves the
product criterion through the two source event arguments.
\end{proof}
\end{annexitem}

\begin{annexitem}{annex:p15.survival-product}{thm:p15-survival-product}{Successive conditioning and the zero-survival case}
\begin{proof}
The infinite survival event is contained in every \(D_n\).
Any zero finite survival probability therefore proves its probability
zero without further conditioning.
Otherwise every \(P(D_n)>0\), including \(P(D_1)=P(A_1^c)>0\).
The scalar definition of conditional probability gives
\[
 P(D_n)=P(D_{n-1})P(A_n^c\mid D_{n-1})
                         =P(D_{n-1})(1-h_n).
\]
Positivity of both survival probabilities forces \(h_n<1\).
Iteration gives the finite product. The events \(D_n\) decrease to
the infinite survival event, so continuity gives the same formula
with the limiting product. Apply the preceding product criterion
to the sequence \(h_2,h_3,\ldots\).
No step factors the probabilities of the original events independently;
the dependence is retained in their successive conditional hazards.
\end{proof}
\end{annexitem}

\begin{annexitem}{annex:p15.partition-ratio-supermartingale}{thm:p15-partition-ratio-supermartingale}{The missing child mass under refinement}
\begin{proof}
First, for integrable \(W\), the elementary average is integrable
because the sum over positive cells of \(\left|\int_IW\,d\nu\right|\)
is at most \(\int|W|\,d\nu\). On each positive cell it integrates
to \(\int_IW\,d\nu\); on a null cell both integrals vanish.
Finite addition gives the all-event identity for the finite
partition sigma-field. Thus the stated elementary conditional
averages require no measure representation theorem.

For the ratio process,
\[
 E_\nu X_n=\sum_{\substack{I\in\mathcal P_n\\\nu(I)>0}}\mu(I)
                                                \leq\mu(\Omega).
\]
Fix a positive cell \(I\in\mathcal P_n\). The cells
\(J\in\mathcal P_{n+1}\) contained in \(I\) partition it, and
\[
 \frac{\int_I X_{n+1}\,d\nu}{\nu(I)}
    =\frac{\sum_{J\subseteq I,\ \nu(J)>0}\mu(J)}{\nu(I)}
    \leq\frac{\mu(I)}{\nu(I)}=X_n|_I.
\]
If local absolute continuity holds at time \(n+1\), every omitted
null child also has zero \(\mu\)-mass, giving equality.
Local absolute continuity at time \(n\) also restores all omitted
terms in the total integral, giving \(\mu(\Omega)\).
Null parent cells need no conditional assertion beyond the chosen
almost-sure version.

In the example, the initial ratio is the common total mass one.
Every positive later cell is \((0,1]\), whose \(\mu\)-mass is zero,
while the cell carrying all target mass is \(\nu\)-null and assigned
ratio zero. Therefore the later ratios vanish and the decrease is strict.
\end{proof}
\end{annexitem}

\begin{annexitem}{annex:p15.elementary-partition-convergence}{thm:p15-elementary-partition-convergence}{An upcrossing proof using only elementary means}
\begin{proof}
If \(M=0\) the assertion is immediate. Otherwise choose rational
\(0\leq a<b\leq M\). Buy one unit of the process when its value is
at most \(a\), hold until its value is at least \(b\), then repeat.
For the step from \(j-1\) to \(j\), let \(H_j\in\{0,1\}\)
indicate whether a unit is held. The decision depends only on
\(Y_0,\ldots,Y_{j-1}\), so \(H_j\) is constant on the finite
\(\mathcal F_{j-1}\) cells. Let \(U_N(a,b)\) be the number of
completed upcrossings by time \(N\), and
\(G_N=\sum_{j=1}^NH_j(Y_j-Y_{j-1})\).
The elementary mean identity on each cell gives
\(E_\nu[H_j(Y_j-Y_{j-1})]=0\), hence \(E_\nu G_N=0\).

Each completed trade gains at least \(b-a\).
An unfinished final trade has gain at least \(Y_N-a\), and thus
at least \(-(a-Y_N)^+\). Consequently, pathwise outside the common
countable null event for the bounds,
\[
              (b-a)U_N(a,b)\leq G_N+(a-Y_N)^+.
\]
Taking expectations gives \(E_\nu U_N(a,b)\leq M/(b-a)\).
Monotone convergence gives the same bound for the total number of
upcrossings, so that number is finite almost surely.
There are countably many rational pairs. On their common
probability-one event, if liminf and limsup differed, rationals
strictly between them would force infinitely many upcrossings,
a contradiction. Thus \(Y_n\) converges to a value in \([0,M]\).
Use the measurable limsup on this event and assign zero on its
complement to define \(Y\). Dominated convergence, with bound \(M\)
for \(|Y_n-Y|\), proves \(L^1\) convergence.
Every expectation calculation above used finite-cell integration;
no existence theorem for general conditional expectation was used.
\end{proof}
\end{annexitem}

\begin{annexitem}{annex:p15.countably-generated-radon-nikodym}{thm:p15-countably-generated-radon-nikodym}{Finite generating partitions, bounded ratios and sigma-finite assembly}
\begin{proof}
First suppose \(\mu\) is finite, \(\nu\) is a probability and
\(\mu\ll\nu\). Put \(c=\mu(\Omega)\),
\(\rho=(\mu+\nu)/(c+1)\).
Let \(\mathcal P_n\) be the finite partition into nonempty Boolean
atoms of \(E_1,\ldots,E_n\), with
\(\mathcal P_0=\{\Omega\}\).
These partitions refine and their sigma-fields generate \(\mathcal F\).

On every positive-\(\rho\) cell \(I\), define the elementary ratios
\(Y_n=\mu(I)/\rho(I)\), \(Z_n=\nu(I)/\rho(I)\).
Assign zero on null cells; their values there are immaterial.
Both ratios lie in \([0,c+1]\), and their sum is \(c+1\)
\(\rho\)-almost surely.
A \(\rho\)-null cell has zero mass under both numerator measures.
The refinement computation of the preceding theorem, with base \(\rho\),
therefore gives exact elementary martingale identities for \(Y_n,Z_n\).
The elementary bounded-convergence lemma gives their bounded measurable
limits \(Y,Z\), almost surely and in \(L^1(\rho)\), with
\(Y+Z=c+1\), \(\rho\)-almost surely.

If \(A\) belongs to the partition sigma-field at time \(m\), then for
\(n\geq m\), finite addition gives \(\int_A Y_n\,d\rho=\mu(A)\).
Passing to the limit gives \(\int_A Y\,d\rho=\mu(A)\).
The union of the increasing finite sigma-fields is a generating algebra;
uniqueness of finite measures on a generating pi-system gives
\(\mu=Y\rho\) on all of \(\mathcal F\).
The identical argument gives \(\nu=Z\rho\).
These are consequences of the finite ratios; they were not introduced
by a Radon--Nikodym theorem.

Now \(\nu(Z=0)=\int_{\{Z=0\}}Z\,d\rho=0\).
Absolute continuity gives \(\mu(Z=0)=0\), hence
\(\rho(Z=0)=0\).
Define \(f=Y/Z\) on \(Z>0\) and zero elsewhere.
It is nonnegative, measurable and finite almost everywhere.
The identity \(\nu=Z\rho\) gives, first for nonnegative simple functions
and then by monotone convergence for every nonnegative measurable \(h\),
\(\int h\,d\nu=\int hZ\,d\rho\).
Therefore
\[
                \int_A f\,d\nu
                 =\int_A fZ\,d\rho
                 =\int_A Y\,d\rho=\mu(A).
\]
This includes \(c=0\) without division by \(c\).
It also identifies \(f\) as the limit of the original partition ratios
\(X_n=\mu(I)/\nu(I)\): local absolute continuity makes every
positive-\(\rho\) cell positive under \(\nu\), so \(X_n=Y_n/Z_n\)
on one common full-\(\rho\) event. Since \(Z>0\) there almost surely,
their quotient converges to \(f\).

For sigma-finite \(\mu,\nu\), take respective countable covers by sets
of finite mass. Enumerate all intersections of one set from each cover
and disjointify them to obtain measurable \(D_i\) partitioning \(\Omega\),
with both \(\mu(D_i),\nu(D_i)<\infty\).
On \(D_i\), the sets \(E_j\cap D_i\) generate the restricted sigma-field.
If \(\nu(D_i)=0\), absolute continuity gives \(\mu(D_i)=0\);
put the local density zero.
Otherwise let \(\lambda_i=\nu|_{D_i}/\nu(D_i)\).
The finite/probability construction gives a finite-almost-everywhere
\(f_i\) representing \(\mu|_{D_i}\) relative to \(\lambda_i\).
Set \(f=f_i/\nu(D_i)\) on this block.
This countable piecewise function is measurable. Summing its
nonnegative integrals over the disjoint blocks yields, for every \(A\),
\[
             \int_A f\,d\nu=\sum_i\mu(A\cap D_i)=\mu(A).
\]
It is finite \(\nu\)-almost everywhere because each local construction
has that property and the exceptional sets form a countable union.

To prove uniqueness, let \(h\geq0\) be another representing function.
On each block its integral is the finite number \(\mu(D_i)\);
therefore \(h\) is finite \(\nu\)-almost everywhere there.
Let \(N\) be the common \(\nu\)-null set where either \(f\) or \(h\)
is nonfinite.
For integers \(k\geq1\), the set
\(A_{i,k}=(D_i\setminus N)\cap\{f\geq h+1/k\}\)
has finite integrals of \(f,h\),
and equality of their represented measures gives
\[
          0=\int_{A_{i,k}}(f-h)\,d\nu\geq\nu(A_{i,k})/k.
\]
Their countable union contains \(\{f>h\}\) apart from the null
nonfinite sets. Interchanging \(f,h\) proves equality almost everywhere.
This completes the independent finite-partition source route.
\end{proof}
\end{annexitem}

\begin{annexitem}{annex:p15.bernoulli-product-complete}{thm:p15-bernoulli-product-complete}{Local support and the exact affinity comparison}
\begin{proof}
If \(\mu\ll\nu\), every coordinate pushforward has the same
absolute-continuity direction: a null coordinate event pulls back to
a \(\nu\)-null cylinder.
For Bernoulli coordinates this says precisely that a zero reference
success probability requires zero target success probability, and a
unit reference success probability requires a unit target probability.
Conversely those implications give local absolute continuity at
every coordinate. With it, the square-root density integral is exactly
\[
       \sqrt{\alpha_m\beta_m}
                     +\sqrt{(1-\alpha_m)(1-\beta_m)}=a_m.
\]
At a reference-null point the integral contribution is zero; the
target mass there is also zero by the support condition.
The general product theorem
\cref{thm:p15-product-affinity-criterion} now gives both directions
of the full criterion. Under local absolute continuity the affinity
is positive, and \cref{thm:p15-product-affinity-series} gives its
equivalent sum condition.

For probabilities \(\alpha,\beta\) in \([\epsilon,1-\epsilon]\),
\[
 2(1-a)
   =(\sqrt\alpha-\sqrt\beta)^2
       +(\sqrt{1-\alpha}-\sqrt{1-\beta})^2
   =(\alpha-\beta)^2
       \left(\frac1{(\sqrt\alpha+\sqrt\beta)^2}
        +\frac1{(\sqrt{1-\alpha}+\sqrt{1-\beta})^2}\right).
\]
Each denominator lies between \(4\epsilon\) and \(4(1-\epsilon)\).
The sum in parentheses therefore lies between
\(1/(2(1-\epsilon))\) and \(1/(2\epsilon)\).
Division by two proves the stated quadratic bounds.
All coordinate laws in this regime are equivalent, so the
product-affinity theorem gives equivalence in the summable case
and singularity in the divergent case.
\end{proof}
\end{annexitem}

\begin{annexitem}{annex:p15.bernoulli-sparse-singularity}{thm:p15-bernoulli-sparse-singularity}{A directly calculated separating tail event}
\begin{proof}
Let \(F=\{\xi_m=1\text{ for only finitely many }m\}\).
It is a measurable tail event.
Under \(\mu\), the first Borel--Cantelli lemma gives \(\mu(F)=1\)
because the success probabilities sum to a finite number.
Under \(\nu\), independence and the second Borel--Cantelli lemma
give \(\nu(F)=0\) because its success probabilities have divergent sum.
Thus the measures are concentrated on disjoint measurable sets,
proving singularity without an affinity approximation.

For the explicit example the target success sum is
\(\sum_{m\geq1}(m+1)^{-2}<\infty\), while the reference success sum
is the divergent harmonic tail. Every coordinate probability lies
strictly between zero and one.
Also \(0<\beta_m-\alpha_m<\beta_m\), so
\[
                  \sum_m(\alpha_m-\beta_m)^2
                       \leq\sum_m\beta_m^2<\infty.
\]
The just-proved singularity therefore disproves a quadratic criterion
without a common interior bound.
\end{proof}
\end{annexitem}

\begin{annexitem}{annex:p15.bernoulli-l1-comparison}{thm:p15-bernoulli-l1-comparison}{Summability without a denominator bound}
\begin{proof}
For nonnegative \(x,y\),
\((\sqrt x-\sqrt y)^2\leq|x-y|\).
Indeed, if \(x\geq y\), then
\(\sqrt x-\sqrt y=(x-y)/(\sqrt x+\sqrt y)\) when the denominator
is positive, and its square is at most \(x-y\); the both-zero
case is immediate. Interchange \(x,y\) for the other order.
Apply this to \(\alpha_m,\beta_m\) and to their complements:
\[
 2(1-a_m)
   =(\sqrt{\alpha_m}-\sqrt{\beta_m})^2
    +(\sqrt{1-\alpha_m}-\sqrt{1-\beta_m})^2
                       \leq2|\alpha_m-\beta_m|.
\]
Thus the affinity deviations have finite sum.
Under the local support condition each affinity is positive,
so the complete criterion gives \(\mu\ll\nu\).
When all probabilities are interior, the reverse local condition
also holds, giving mutual absolute continuity.

For the final example \(a_1=1/\sqrt2\) and all later affinities
are one. The cylinder \(\{\xi_1=1\}\) has \(\mu\)-mass \(1/2\)
and \(\nu\)-mass zero, disproving \(\mu\ll\nu\).
They are not singular: the common nonzero measure
\(\eta=\tfrac12\nu\) is dominated by both.
To see the second domination, \(\mu\) is the half-and-half mixture
of the laws with first coordinate fixed at zero or at one, and the
first component is \(\tfrac12\nu\).
Two singular measures cannot dominate the same nonzero measure:
a separating support would force that common measure to vanish
both on it and on its complement.
This establishes both claimed failures.
\end{proof}
\end{annexitem}
